\documentclass[10pt,a4paper]{amsart}
\usepackage[margin=19mm]{geometry}
\usepackage{amsmath,amssymb,mathtools}
\usepackage[scr=rsfs,cal=euler]{mathalfa}
\usepackage{xfrac}
\usepackage[unicode,hidelinks,bookmarks=false]{hyperref}
\newcommand{\ie}{i.e.\ }
\newcommand{\cf}{cf.\ }
\newcommand{\resp}{respectively\ }
\newcommand{\Subsection}{\subsection}
\providecommand{\nobreakdash}{\nobreak\hbox{-}}
\setlength{\emergencystretch}{2em}
\multlinegap0pt
\usepackage{mathtools}
\usepackage{float}
\usepackage{xcolor}
\usepackage{amssymb}
\usepackage[shortlabels]{enumitem}
\newtheorem{theorem}{Theorem}
\title{SIGMA: Scalable Spectral Insights for LLM Model Collapse}
\author{Yi Gu, Lingyou Pang, Xiangkun Ye, Tianyu Wang, Jianyu Lin, Carey E. Priebe, Alexander Aue}
\date{Source: arXiv:2601.03385v3 (CC BY 4.0)}
\begin{document}
\maketitle

\noindent\emph{Source and licence.} SIGMA: Scalable Spectral Insights for LLM Model Collapse. Version arXiv:2601.03385v3 (CC BY 4.0), distributed by arXiv under the Creative Commons Attribution 4.0 International licence, http://creativecommons.org/licenses/by/4.0/, as recorded in the arXiv licence metadata for this exact version and shown on the abstract page https://arxiv.org/abs/2601.03385v3. Full licence notice: \texttt{licenses/arxiv-2601.03385-CC-BY-4.0.txt}.\par\smallskip

\noindent\textit{Editorial scope.} Counted unit: the article's theorem thm:stochastic (source0.tex lines 219--233). The article's complete original proof of it is delivered with it (source0.tex lines 621--657, one \texttt{proof} environment of 37 lines, which the article places away from the statement). No mathematics is written, repaired or re-derived: the statement and the proof are the article's own lines, in the article's own notation, and no retained line is substituted or re-flowed.  No further source-local context is needed: every cross-reference of every retained line resolves inside the two delivered spans.  Nothing was omitted from any retained span, so the package carries no omission record.  No retained line contains a citation, so the wrapper carries no bibliography and no bibliography entry.  No retained line contains a figure environment, an image-inclusion command or drawing code, so no image is delivered and the package declares no asset dependency.  Editorial formatting only, in addition: an AMS \texttt{amsart} preamble that restates the article's own declarations and macros used by the retained lines, namely the environments \texttt{theorem} on separate counters, together with the standard packages those lines need.  The retained environments print in this document as Theorem 1 in order of appearance, whereas the article prints the same items with the numbering of its own full document.  The complete native source of the article as distributed in the arXiv e-print for this exact version is bundled unchanged as \texttt{sources/192222/source0.tex}.
\begin{theorem}[Stochastic Scaling Spectral Bound]
\label{thm:stochastic}
Assume the columns of $M^{(k)}$ are i.i.d. samples drawn from a distribution with positive-definite covariance matrix $C \in \mathbb{S}^m_{+}$. As $n_A, n_k \to \infty$ with fixed $m$, for any $\delta \in (0, 1)$, there exists a constant $K$ such that with probability at least $1-\delta$:
\begin{align}
    \det(G^{(k)}) \le K \cdot \left(\frac{n_k}{n_A}\right)^m \det(G_A).
\end{align}
% there exists a finite constant $K$ such that with high probability:
% \begin{align}
%     \det(G^{(k)}) \ \approx \ \left( \frac{n_k}{n_A} \right)^m \det(G_A^{(k)}).
% \end{align}
% Specifically, the upper bound holds in expectation:
% \begin{align}
%     \mathbb{E}[\det(G^{(k)})] \ \leq \ K \cdot \left(\frac{n_k}{n_A} \right)^m \det(G_A^{(k)}).
% \end{align}
\end{theorem}

\begin{proof}[Proof of Theorem \ref{thm:stochastic}]
We leverage the concentration of the sample covariance matrix. Let $\hat{C}_{n_k} = \frac{1}{n_k} G^{(k)}$ and $\hat{C}_{n_A} = \frac{1}{n_A} G^{(k)}_A$. By the properties of determinant scaling, we have:
\begin{align}
    \det(G^{(k)}) &= \det(n_k \hat{C}_{n_k}) = n_k^m \det(\hat{C}_{n_k}), \nonumber \\
    \det(G^{(k)}_A) &= \det(n_A \hat{C}_{n_A}) = n_A^m \det(\hat{C}_{n_A}).
\end{align}
We consider the ratio:
\begin{align}
    \frac{\det(G^{(k)})}{\det(G^{(k)}_A)} = \left( \frac{n_k}{n_A} \right)^m \frac{\det(\hat{C}_{n_k})}{\det(\hat{C}_{n_A})}.
\end{align}
By the Strong Law of Large Numbers (or finite-sample concentration inequalities for covariance matrices), both $\hat{C}_{n_k}$ and $\hat{C}_{n_A}$ converge to the population covariance $C$ as $n_k, n_A \to \infty$.
Specifically, the continuous mapping theorem implies that $\det(\hat{C}_{n_k}) \xrightarrow{\mathbb{P}} \det(C)$ and $\det(\hat{C}_{n_A}) \xrightarrow{\mathbb{P}} \det(C)$. This implies that for any $\epsilon > 0$, with probability approaching 1:
\begin{align}
    1 - \epsilon \le \frac{\det(\hat{C}_{n_k})}{\det(\hat{C}_{n_A})} \le 1 + \epsilon.
\end{align}
Letting $K = 1+\epsilon$ establishes the upper bound with high probability:
\begin{align}
    \det(G^{(k)}) \le K \left( \frac{n_k}{n_A} \right)^m \det(G^{(k)}_A).
\end{align}
This confirms that the determinant of the full Gram matrix scales with the determinant of the sub-sample according to the geometric factor $(n_k/n_A)^m$, up to a stochastic constant $K$ that approaches 1.
% By the Strong Law of Large Numbers,
% \begin{align}
%     \frac{1}{n_k}G^{(k)} \to C \quad \text{a.s.}
% \end{align}
% For the original Gram matrix, since $n_k$ is usually very large, we have
% \begin{align}
%     \det(G^{(k)}) \approx \det(n_kC) = n_k^m \det(C).
% \end{align}
% On the other hand, assuming $n_A$ is sufficiently large to approximate the population covariance, then
% \begin{align}
%     \det(G_A^{(k)}) \approx O(\det(n_AC)) = n_A^m \det(C).
% \end{align}
% Therefore, we have a stochastic upper bound of the original Gram matrix using the sub-Gram matrix: there exists a finite constant $K < \infty$ such that for each iteration $k$,
% \begin{align}
%     \det(G^{(k)}) \leq K \cdot \left(\frac{n_k}{n_A} \right)^m \det(G_A^{(k)}).
% \end{align}
\end{proof}

\end{document}
